% HD-R07. GPT-6 Astra Pro, /root; inherited effort; 21 November 2026.
\subsection{HD-R07: every character polynomial in one ample sixfold class}
\label{sec:ample-polynomial-sixfold-orbits}

Let $F$ be an actual perfect complex on a complex abelian sixfold
whose whole character is
\[
                 \operatorname{ch}(F)=e^{B_0}P_0(H),           \tag{R07.1}
\]
where $B_0$ is a real $(1,1)$ class, $H$ is an ample real
$(1,1)$ class, and $P_0\in\mathbf R[z]$ has degree at most six.
The displayed whole character is rational because it is the
character of the actual complex. No finite composition of actual
Fourier--Mukai equivalences, isogeny pullbacks or pushforwards,
arbitrary line-bundle twists, translations, shifts and derived
duals carries it, even up to nonzero scalar, to
\[
 e^B\bigl(P(\theta)+\lambda t\bigr),\quad
 \deg P\le6,\quad0\ne\lambda t\in W_K(A),                    \tag{R07.2}
\]
on a polarized Weil sixfold of signature $(3,3)$.
This extends the ample imaginary-secant source test HD-N18 to
every whole character (R07.1). It supplies no assertion about
other characters or correspondences outside the specified orbit.

\paragraph{1. The exact source Hochschild map.}
Undo the displayed real $(1,1)$ shear using (N18.4). Choose
bases $x_i\in H^{1,0}$, $y_i\in H^{0,1}$ with
$H=\sum_{i=1}^6x_i y_i$. The twelve degree-one generators
are the wedges $y_i\wedge$ and contractions $\iota_{x_i}$.
Their degree-two products act on $P_0(H)$ by
\[
 y_i y_j P_0(H),\quad
 \pm y_i y_j P_0'(H),\quad
 \pm y_i y_j P_0''(H)\qquad(i<j),                            \tag{R07.3}
\]
with the middle column coming from the two ordered mixed
operators. A mixed operator with equal indices gives zero.
These are actual exterior differentiation identities: contraction
of $H$ is $y_i$, and contraction of a $y_j$ is zero.

For $0\le k\le4$, the map
$\eta\mapsto\eta H^k$ on $\bigwedge^2\langle y_i\rangle$
is injective. Given the coefficient of $y_i y_j$, choose $k$
paired indices avoiding $i,j$ and read the monomial containing
exactly those holomorphic factors. Its coefficient recovers the
chosen coefficient of $\eta$, multiplied by the nonzero $k!$.
Different $k$ have different bidegrees. Thus the full image in
(R07.3) is the tensor product of this fifteen-dimensional
coefficient space with the span of the three truncated columns
$P_0,P_0',P_0''$ modulo $z^5$.

Write $P_0(z)=\sum_{j=0}^6p_jz^j$ and $m_j=j!p_j$. Row
scaling by $k!$ identifies these columns with the actual matrix
\[
 M_0=(m_{k+j})_{0\le k\le4,\ 0\le j\le2},\qquad
           \operatorname{rank}c_{P_0(H)}^2
                       =15\operatorname{rank}M_0.            \tag{R07.4}
\]
No injectivity of a self-Ext trace is assumed here; this is the
explicit Hochschild action on the whole character.

For (R07.2), HD-R06 gives the exact rank set
$\{30,39,48,57\}$. The only member also in
$\{0,15,30,45\}$ is $30$. The actual Hodge Clifford maps of
HD-N18 preserve this rank, including the duality sign. Therefore
a proposed transport forces
\[
             \operatorname{rank}M_0=2,
             \qquad P=0\text{ in (R07.2)}.                 \tag{R07.5}
\]
The complete target character, after its specified shear, is
therefore a nonzero rational pure Weil class. Its full real spin
stabilizer is $\mathfrak{su}(6,6)$ with Killing signature
$(72,71)$, through the literal Fock map and hyperbolic basis of
HD-R06. We use that proved map, rather than an identification
based on its dimension.

\paragraph{2. The actual recurrence exhausts the source characters.}
Choose a nonzero real kernel vector $(q_0,q_1,q_2)$ of $M_0$.
It gives exactly five equations
\[
             q_0m_k+q_1m_{k+1}+q_2m_{k+2}=0,
                                  \qquad0\le k\le4.         \tag{R07.6}
\]
We solve these finite equations, retaining every endpoint case.
All exponentials below mean their finite exterior expansions
modulo $H^7=0$.

With $q_2\ne0$, the first two moments determine all seven by
the displayed second-order recurrence. Let $r_1,r_2$ be the
roots of $q_2r^2+q_1r+q_0$. For distinct roots, solving the
two initial equations gives
\[
 P_0(H)=\alpha e^{r_1H}+\beta e^{r_2H}.                       \tag{R07.7}
\]
Both coefficients are nonzero, since a single exponential has
matrix rank at most one. For a repeated nonzero root $r$, the
recurrence solutions are $m_k=ar^k+bk r^{k-1}$. Their exponential
generating function is
\[
                         P_0(H)=(a+bH)e^{rH},\quad b\ne0.    \tag{R07.8}
\]
This last expression also covers $r=0$ directly: its moments are
$m_0=a,m_1=b,m_k=0$ for $k\ge2$. Substitution verifies the
five equations in that case, without division by $r$.

With $q_2=0,q_1\ne0$, put $r=-q_0/q_1$. The equations give
$m_k=m_0r^k$ for $0\le k\le5$, with $m_6$ free. Hence
\[
                         P_0(H)=a e^{rH}+bH^6.               \tag{R07.9}
\]
Both $a,b$ are nonzero under rank two; otherwise its Hankel
matrix has rank at most one. Finally $q_2=q_1=0$ forces
$m_0=\cdots=m_4=0$, and gives
\[
                         P_0(H)=aH^5+bH^6,\quad a\ne0.       \tag{R07.10}
\]
These cases exhaust the chosen nonzero kernel vector. They are
an exact finite recurrence calculation; no classification of
secant varieties is invoked.

\paragraph{3. Distinct exponential roots and the real pair cases.}
The coefficients of (R07.6) are real. Its distinct roots are
therefore either both real or a conjugate pair $r\pm ic$ with
$c>0$. In the conjugate case, (R07.7), with $e^{rH}$ absorbed
by the explicit real $(1,1)$ shear, is precisely the class of
HD-N18. That whole proof, including the recovered pure axes and
the incompatible definite and neutral symmetric forms, excludes
the proposed transport.

For two distinct real roots, their pure-spinor annihilators are
the actual real graphs
\[
             L_j=\{-r_jH^\sharp\xi+\xi:\xi\in V^*\},
                    \qquad V=H^1(A,\mathbf R).
\]
They are transverse because $(r_1-r_2)H^\sharp$ is invertible,
and are maximal isotropic by skewness of $H^\sharp$.
The map $\bigwedge^*L_2\to\bigwedge^*V$ obtained by applying
its Clifford operators to $e^{r_1H}$ is the actual relative Fock
map of HD-N17. After conjugating by $e^{-r_1H}$, its highest
exterior terms are the exterior powers of the invertible map
$(r_2-r_1)H^\sharp$, up to the fixed overall sign. Hence it
is an isomorphism over $\mathbf R$. It sends the two source
lines to the vacuum and top lines, with both coefficients nonzero.
Their vector stabilizer is exactly
$\mathfrak{sl}_{12}(\mathbf R)$, by the separate degree-two,
degree-ten and trace equations proved in HD-R06. Its Killing
signature $(77,66)$ differs from the target signature $(72,71)$.
The actual functor maps conjugate the full stabilizers, so this
case is excluded.

For (R07.9), the real $(1,1)$ shear $e^{-rH}$ gives
$a+bH^6$, since $e^{-rH}H^6=H^6$. The two real annihilators
are respectively $V^*$ and $V$, already transverse. The same
vacuum--top calculation gives exactly
$\mathfrak{sl}_{12}(\mathbf R)$ and the same exclusion.

\paragraph{4. Exact maps for both tangent cases.}
Take a real Darboux basis $H=\sum_{i=1}^6a_i b_i$, and define
the actual operators
\[
 L_i=(a_i b_i)\wedge,\qquad
 \Lambda_i=\iota_{b_i^*}\iota_{a_i^*},\qquad
 L=\sum_iL_i,\quad\Lambda=\sum_i\Lambda_i.
\]
These are respectively wedge and contraction bivectors in the
full spin Lie algebra. Their finite exponentials are invertible
spin operators, so they give explicit conjugations of the Lie
action and its stabilizers. We use them to compute a stabilizer;
we make no assertion that each is the transform of a geometric
equivalence.

After removing $e^{rH}$ in (R07.8), the actual operator
$\exp(-a\Lambda/(6b))$ carries $a+bH$ to $bH$, because
$\Lambda H=6$ and a second contraction kills $H$.
For (R07.10), use the explicit invertible operator
\[
                 \mathcal F_H=e^L e^{-\Lambda}e^L.
\]
On the two-dimensional span $\langle1,a_i b_i\rangle$ in
each pair it sends $1$ to $a_i b_i$ and $a_i b_i$ to $-1$.
The six even pair operators commute, and on their product this
gives the exact identity
\[
 \mathcal F_H(e^{tH})=t^6 e^{-H/t},\qquad
 \mathcal F_H(H^k)=(-1)^{6-k}\frac{k!}{(6-k)!}H^{6-k}.
                                                               \tag{R07.11}
\]
This is a polynomial identity after multiplication by $t^6$;
it does not require assigning a value at $t=0$ to a denominator.
In particular (R07.10) goes to $-120aH+720b$. The previous
explicit contraction removes its constant term. Both tangent
cases therefore have the actual spin stabilizer of $H$, through
specified invertible real operators.

\paragraph{5. The full stabilizer of the symplectic two-form.}
Write a general spin Lie element on $\bigwedge^*V$ as
$\beta\wedge+D_A-\operatorname{tr}(A)/2+\iota_\gamma$.
Its three outputs on $H$ have distinct exterior degrees four,
two and zero. The map $\beta\mapsto\beta H$ is injective.
Indeed the explicit Darboux operators give
$[\Lambda,L]=6-k$ on exterior degree $k$. For $\beta$ of
degree two with $L\beta=0$, this gives
$4\beta=-L\Lambda\beta$. Since $\Lambda\beta$ is a scalar,
wedging once more and using $H^2\ne0$ forces that scalar and
then $\beta$ to vanish. Thus every stabilizer element has
$\beta=0$.

The degree-two equation is
$D_AH=\operatorname{tr}(A)H/2$. For a conformally symplectic
endomorphism satisfying $D_AH=kH$, contraction of its matrix
equation with the inverse form gives $2\operatorname{tr}A=12k$.
Combining the two equations gives $\operatorname{tr}A=0$ and
$k=0$. Hence $A\in\mathfrak{sp}(V,H)$. The degree-zero
equation is exactly $\iota_\gamma H=0$. Consequently the full
real stabilizer is the literal semidirect product
\[
 \mathfrak s_H=\mathfrak{sp}_{12}(\mathbf R)
          \ltimes\{\gamma\in\bigwedge^2V^*:\iota_\gamma H=0\}.
                                                               \tag{R07.12}
\]
The first summand has dimension $36+21+21=78$ in its usual
symplectic block matrices. The second has dimension
$\binom{12}{2}-1=65$; it is nonzero and abelian since the even
contraction operators commute. The derivation commutator preserves
it, so it is an actual proper Lie ideal. The total dimension is
$143$.

The target pure-Weil stabilizer $\mathfrak{su}(6,6)$ also has
dimension $143$ and is simple: its explicit complexification is
$\mathfrak{sl}_{12}(\mathbf C)$, as in HD-N17's
Hermitian-adjoint scalar-extension formula. Conjugation through
the actual transport and the explicit operators above would be
a Lie isomorphism, carrying the proper nonzero abelian ideal in
(R07.12) to such an ideal in this simple algebra. The latter has
none. This excludes both tangent cases and completes the proof.

\paragraph{Checks and exact implication.}
Every recurrence, source case and comparison above uses an explicit
map. The target ranks and real stabilizer are the complete
sixfold results HD-R06; the Hodge-compatible transport and imaginary
pair obstruction are HD-N18. The latter in turn uses the original
Fourier graph maps of GLO Proposition4.3.7 and Orlov Theorem2.10
at the exact locators already given in HD-G19. This result closes
every proposed source with whole character polynomial in one ample
class, including any actual bundle or monad that has that character.
It does not decide the existence of such bundles and does not
extend the assertion to an arbitrary polynomial in several divisor
classes.
